% ..............................................................................
% Demo of the fau-beamer template.
%
% Copyright 2022 by Tim Roith <tim.roith@fau.de>
%
% This program can be redistributed and/or modified under the terms
% of the GNU Public License, version 2.
%
% ------------------------------------------------------------------------------

% ------------------------------------------------------------------------------
\section{Example Section}
% ..............................................................................
\begin{frame}[t]{An Example Slide}{With a subtitle}
This is how a normal slide looks like, where plain text is put within the \texttt{frame} environment.
\end{frame}
%
%
%
%
%
% ..............................................................................
\begin{frame}[t]{Blocks}{}
You can use blocks with this template.
% ----------------
\begin{block}{Block headline}
This block spans all the way from the left to the right margin.
\end{block}
% ..............................................................................
Using minipages one can use blocks of certain fixed sizes.\\
%
\noindent%
\begin{minipage}[t]{0.45\textwidth}
\begin{block}{Headline A}
Some text.
\end{block}
\end{minipage}
\hfill%
% --------
\begin{minipage}[t]{0.2\textwidth}
\begin{block}{Headline B}
Some other text\\
that spans\\
over multiple lines.
\end{block}
\end{minipage}
\hfill%
% --------
\begin{minipage}[t]{0.2\textwidth}
\begin{block}{Headline C}
Even more text\\
that spans\\
over multiple lines.
\end{block}
\end{minipage}

\begin{block}{}
A block with an empty title.
\end{block}

\end{frame}
%
%
%
%
%
% ..............................................................................
\begin{frame}{A very very very very very very very long title}

\end{frame}
%
%
%
%
%
% ..............................................................................
\begin{frame}[t]{Top aligned}
In this frame everything is aligned on top.
\end{frame}
%
%
%
%
%
% ..............................................................................
\begin{frame}[b]{Bottom aligned}
In this frame everything is aligned on the bottom.
\end{frame}
%
%
%
%
%
% ..............................................................................
\begin{frame}{Lists}{Itemize and enumerate}
You can use the itemize environment that looks as follows.
% --------------------------------------
\begin{itemize}
\item An item.
\item Another one.
\begin{itemize}
\item A subitem.
\item Another subitem.
\end{itemize}
\item And another item.
\end{itemize}
%
\vfill%
You can also use the enumerate environment.
% --------------------------------------
\begin{enumerate}
\item The first item.
\item The second one.
\item The third one.
\end{enumerate}
% --------------------------------------
\end{frame}

% ..............................................................................
\begin{frame}{Using Overlay Specifications}
\begin{itemize}[<+(1)->]
\item The first item.
\item The second one.
\end{itemize}
\pause
\begin{block}{Framed Text}
This should be displayed after the list.
\pause
This should be displayed last.
\end{block}
\end{frame}
%
%
%
%
%
% ..............................................................................
\subsection{Subsection A}
\begin{frame}{\mbox{}}
We can also have subsections, even though they don't really do anything right now. They only appear in the toc. 
\end{frame}
%
%
%
%
%
% ..............................................................................
\subsection{Subsection B}
\begin{frame}{\mbox{}}
When there is one subsection, there probably should be another.
\end{frame}
%
%
\begin{frame}{Graphics}
This slide demonstrates graphics and captions.
\begin{figure}
\includegraphics[width=.7\textwidth]{template-art/FAUWortmarkeBlau.pdf}
\caption{The FAUWortmarkeBlau file.}
\end{figure}
\end{frame}
%
%
%
%
%
% ------------------------------------------------------------------------------
\section{The Color Scheme}
%
%
%
%
%
% ..............................................................................
\begin{frame}[fragile]{Choosing the Institution}
You can specify the institute template by passing it to package, i.e.,
%
\begin{center}
\begin{verbatim}
\usepackage[institute=<option>]{styles/beamerthemefau},
\end{verbatim}
\end{center}
%
where you have the following options,
\begin{itemize}
\item FAU
\item RW
\item Med
\item Nat
\item TF
\end{itemize}
\end{frame}
%
%
%
%
%
% ..............................................................................
\begin{frame}{The Colors}{Base and Dark scheme}
For each institute the color scheme consists of two main colors, which are named 
\texttt{BaseColor} and \texttt{BaseDarkColor}. They can be used throughout the document. For each of theses colors, adding a letter from A to D will create a lighter shade as displayed below.
%
%
%
%
%
\newdimen\bwidth
\newdimen\bheight
\newdimen\pdim
\setlength{\bwidth}{0.35\textwidth}
\setlength{\bheight}{0.5\textheight}
\setlength{\pdim}{0.05\textwidth}
\vspace{1em}%
\begin{columns}
\begin{column}{0.45\textwidth}
\setbeamercolor{colordemo}{bg=BaseColor}
\begin{beamercolorbox}[dp=0pt,sep=0em,wd=\bwidth,ht=\bheight]{}
%
\begin{beamercolorbox}[sep=0em,wd=\pdim,ht=\pdim]{colordemo}%
\end{beamercolorbox}\quad\texttt{BaseColor}\newline%
\setbeamercolor{colordemo}{bg=BaseColorA}%
\begin{beamercolorbox}[sep=0em,wd=\pdim,ht=\pdim]{colordemo}
\end{beamercolorbox}\quad\texttt{BaseColorA}\hfill\newline%
\setbeamercolor{colordemo}{bg=BaseColorB}%
\begin{beamercolorbox}[sep=0em,wd=\pdim,ht=\pdim]{colordemo}%
\end{beamercolorbox}\quad\texttt{BaseColorB}\hfill\newline%
\setbeamercolor{colordemo}{bg=BaseColorC}%
\begin{beamercolorbox}[sep=0em,wd=\pdim,ht=\pdim]{colordemo}%
\end{beamercolorbox}\quad\texttt{BaseColorC}\hfill\newline%
\setbeamercolor{colordemo}{bg=BaseColorD}%
\begin{beamercolorbox}[sep=0em,wd=\pdim,ht=\pdim]{colordemo}%
\end{beamercolorbox}\quad\texttt{BaseColorD}%
\end{beamercolorbox}%
\end{column}
%
% 
%
\begin{column}{0.45\textwidth}
\setbeamercolor{colordemo}{bg=BaseDarkColor}
\begin{beamercolorbox}[dp=0pt,sep=0em,wd=\bwidth,ht=\bheight]{}
%
\begin{beamercolorbox}[sep=0em,wd=\pdim,ht=\pdim]{colordemo}%
\end{beamercolorbox}\quad\texttt{BaseDarkColor}\newline%
\setbeamercolor{colordemo}{bg=BaseDarkColorA}%
\begin{beamercolorbox}[sep=0em,wd=\pdim,ht=\pdim]{colordemo}
\end{beamercolorbox}\quad\texttt{BaseDarkColorA}\hfill\newline%
\setbeamercolor{colordemo}{bg=BaseDarkColorB}%
\begin{beamercolorbox}[sep=0em,wd=\pdim,ht=\pdim]{colordemo}%
\end{beamercolorbox}\quad\texttt{BaseDarkColorB}\hfill\newline%
\setbeamercolor{colordemo}{bg=BaseDarkColorC}%
\begin{beamercolorbox}[sep=0em,wd=\pdim,ht=\pdim]{colordemo}%
\end{beamercolorbox}\quad\texttt{BaseDarkColorC}\hfill\newline%
\setbeamercolor{colordemo}{bg=BaseDarkColorD}%
\begin{beamercolorbox}[sep=0em,wd=\pdim,ht=\pdim]{colordemo}%
\end{beamercolorbox}\quad\texttt{BaseDarkColorD}%
\end{beamercolorbox}%
\end{column}
\end{columns}
\end{frame}
%
%
%
%
%
% ------------------------------------------------------------------------------
\section{The Frame Dimensions}
%
\setbeamercolor{colordemo}{fg=black,bg=BaseDarkColorC}
%
%
%
\begin{frame}[t]{Text Area}{The area you can fill in a normal frame}

\begin{beamercolorbox}[dp=0pt,sep=5em,wd=\textwidth,ht=\FrameHeight,center]{colordemo}%
This area can be filled on a normal frame.
\end{beamercolorbox}%
\end{frame}
%
%
%
\begin{frame}

\begin{beamercolorbox}[dp=0pt,sep=5em,
					   wd=\textwidth,
					   ht=\dimexpr\FrameHeight+\TitleHeight+4mm,center]{colordemo}%
This area can be filled on a frame without a title.
\end{beamercolorbox}%
\end{frame}
%
%
%
%
%
\begin{frame}[plain]
\begin{beamercolorbox}[dp=0pt,sep=5em,wd=\textwidth,ht=\paperheight,center]{colordemo}%
This area can be filled on a plain frame.\\
In order to achieve this one has to use the option \texttt{[plain]} for the frame.
\end{beamercolorbox}%
\end{frame}
%
%
%
%
%
\begin{frame}[blank]
\begin{beamercolorbox}[dp=0pt,sep=5em,wd=\textwidth,ht=\paperheight,center]{colordemo}%
This area can be filled on a blank frame with removed margins.\\
In order to achieve this one has to use the option \texttt{[blank]} for the frame. This is option 
is specifically defined in the \texttt{fau-beamer} template and doesn't work in the standard 
\texttt{beamer} class.
\end{beamercolorbox}%
\end{frame}
%
%
%
%
%
% ------------------------------------------------------------------------------
\section{Equations and blocks}
%
%
%
%
%
% ..............................................................................
\begin{frame}{Equations}{Label alignment}
The \texttt{beamer} class unfortunately has some problems with the proper 
alignment of equation labels.

%
\pause
%
In this template we try to patch some of these issues and therefore the label 
below should be aligned exactly on the right margin
%
\begin{equation}\label{eq:pythagoras}
a^2 + b^2 = c^2.
\end{equation}
%
An \texttt{align} environment should yield the same alignment
%
\begin{align}\label{eq:euler}
e^{i\pi} +1= 0.
\end{align}

This is especially important if we put such an environment in a block,

\begin{block}{Gauss's Theorem}
\begin{align}\label{eq:gauss}
\int_\Omega \div F\ d\lambda^n = \int_{\partial\Omega} F\cdot \nu\ d\mathcal{H}^{n-1}.
\end{align}
\end{block}
\end{frame}
% ------------------------------------------------------------------------------
\section{Citing and bibliography}
%
%
%
%
%
\begin{frame}{Citing}
Since we use \texttt{biblatex} we can easily cite our favourite articles, like
\cite{bungert2021bregman, bungert2021neural}. For presentation the \texttt{footnote} command is useful to have citation at the bottom\footnote{\cite{bungert2021bregman}}.
\end{frame}
%
%
%
%
%
\begin{frame}{References}
To show the references used within this presentation we can use the \texttt{printbibliography} command from \texttt{biblatex}. For beamer documents it is important to give the additional option \texttt{[heading=none]}.
%
%
%
%
%
\begin{block}{References}
\printbibliography[heading=none]
\end{block}
\end{frame}
%
%
%
%
%
\begin{frame}{fau-beamer}
%
\end{frame}